\documentclass[10pt,a4paper]{amsart}
\usepackage[margin=19mm]{geometry}
\usepackage{amsmath,amssymb,mathtools}
\usepackage[scr=rsfs,cal=euler]{mathalfa}
\usepackage{xfrac}
\usepackage[unicode,hidelinks,bookmarks=false]{hyperref}
\newcommand{\ie}{i.e.\ }
\newcommand{\cf}{cf.\ }
\newcommand{\resp}{respectively\ }
\newcommand{\Subsection}{\subsection}
\providecommand{\nobreakdash}{\nobreak\hbox{-}}
\setlength{\emergencystretch}{2em}
\multlinegap0pt
\usepackage{amssymb}
\usepackage{algorithm}\usepackage{algpseudocode}
\usepackage{xcolor}
\newtheorem{prop}{Proposition}
\title{Szego Theorem for Operator Orthogonal Polynomials}
\author{Badr Missaoui, Nicholas H. Bingham}
\date{Source: arXiv:2305.19025v4 (CC BY 4.0)}
\begin{document}
\maketitle

\noindent\emph{Source and licence.} Szego Theorem for Operator Orthogonal Polynomials. Version arXiv:2305.19025v4 (CC BY 4.0), distributed by arXiv under the Creative Commons Attribution 4.0 International licence, http://creativecommons.org/licenses/by/4.0/, as recorded in the arXiv licence metadata for this exact version and shown on the abstract page https://arxiv.org/abs/2305.19025v4. Full licence notice: \texttt{licenses/arxiv-2305.19025-CC-BY-4.0.txt}.\par\smallskip

\noindent\textit{Editorial scope.} Counted unit: the article's prop prop\_kernelCD (source0.tex lines 567--597). The article's complete original proof of it is delivered with it (source0.tex lines 599--695, one \texttt{proof} environment of 97 lines). No mathematics is written, repaired or re-derived: the statement and the proof are the article's own lines, in the article's own notation, and no retained line is substituted or re-flowed.  Retained uncounted as the source-local context the proof needs: lines 319--336.  Nothing was omitted from any retained span, so the package carries no omission record.  No retained line contains a citation, so the wrapper carries no bibliography and no bibliography entry.  No retained line contains a figure environment, an image-inclusion command or drawing code, so no image is delivered and the package declares no asset dependency.  Editorial formatting only, in addition: an AMS \texttt{amsart} preamble that restates the article's own declarations and macros used by the retained lines, namely the environments \texttt{prop} on separate counters, together with the standard packages those lines need.  The retained environments print in this document as Proposition 1 in order of appearance, whereas the article prints the same items with the numbering of its own full document.  The complete native source of the article as distributed in the arXiv e-print for this exact version is bundled unchanged as \texttt{sources/142242/source0.tex}.
\begin{equation}\label{schurcomplement}
\begin{pmatrix}
A & B\\
B^{*} & C
\end{pmatrix}
=\begin{pmatrix}
I & 0\\
B^{*}A^{-1} & I
\end{pmatrix}
\begin{pmatrix}
A & 0\\
0 & C-B^{*}A^{-1}B
\end{pmatrix}
\begin{pmatrix}
I & A^{-1}B\\
0 & I
\end{pmatrix}.
\end{equation}

\begin{prop}\label{prop_kernelCD}
Define the right and left kernel polynomials of degree $n$ as
$$K_{n}^{R}(w,z) = \sum_{k=0}^{n}\varphi_{k}^{R}(z)\varphi_{k}^{R}(w)^{\dag} ,\quad K_{n}^{L}(w,z) = \sum_{k=0}^{n}\varphi_{k}^{L}(z)^{\dag}\varphi_{k}^{L}(w),$$
Then, we have
\begin{enumerate}
%[label=(\roman*)]
\item
$
K_{n}^{R}(w,z)=[I ~ zI ~\cdots ~ z^{n}I]~T_{n+1}^{-R}~\begin{bmatrix}
I \\
w^{-1}I \\
\vdots \\
w^{-n}I
\end{bmatrix}.
$
\item 
$
K_{n}^{L}(w,z)=[I ~ zI ~\cdots ~ z^{-n}I]~T_{n+1}^{-L}~\begin{bmatrix}
I \\
w^{1}I \\
\vdots \\
w^{n}I
\end{bmatrix}.
$
\item \begin{align} \label{cd_right}
K_{n}^{L}(w,z) &=&\frac{\varphi_{n}^{R,*}(w)^{\dag}\varphi_{n}^{R,*}(z) - \bar{w}z\varphi_{n}^{L}(w)^{\dag}\varphi_{n}^{L}(z)}{1-\bar{w}z}, \\ 
\label{cd_left}
K_{n}^{R}(w,z) &=&\frac{\varphi_{n}^{L,*}(w)\varphi_{n}^{L,*}(z)^{\dag} - \bar{w}z\varphi_{n}^{R}(w)\varphi_{n}^{R}(z)^{\dag}}{1-\bar{w}z}. 
\end{align}
\end{enumerate}
\end{prop}

\begin{proof}
We write $$
K_{n}^{R}(w,z)=[Z ~ z^{n}I]~T_{n+1}^{-R}~\begin{bmatrix}
W^{*}\\
w^{-n}I
\end{bmatrix},
$$ with $Z=[I ~ zI ~\cdots ~ z^{n-1}I]$ and $W=[I ~ wI ~\cdots ~ w^{n-1}I]$.\\
Writing  
$$
T^{R}_{n+1}=
\begin{pmatrix}
T^{R}_{n} & \nu_{n}\\
\nu^{*}_{n} & \mu_{0}
\end{pmatrix}
$$ 
and using (\ref{schurcomplement}), we have 
$$
T^{-R}_{n+1}=
\begin{pmatrix}
A & \gamma\\
\gamma^{*} & \alpha
\end{pmatrix},
$$
where 
$$
A = T^{-R}_{n} + T^{-R}_{n}\nu_{n}\kappa_{n}^{-R}\nu_{n}^{*}T^{-R}_{n},~~~ \gamma = -T^{-R}_{n}\nu_{n}\kappa_{n}^{-R},~~~ \alpha = \kappa_{n}^{-R}.
$$
Using the equality 
$$
[I ~ zI ~\cdots ~ z^{n-1}I]T^{-R}_{n}\nu_{n}\kappa_{n}^{-R/2}
= z^{n}\kappa_{n}^{-R/2} - \varphi_{n}^{R}(z),
$$
we rewrite $K_{n}^{R}(w,z)$ as follows:
\begin{align*}
K_{n}^{R}(w,z)
&= ZAW^{*} + z^{n}\gamma^{*}X^{*}+Z\gamma w^{-n} + z^{n}\alpha x^{-n}\\
&= Z \Bigl(T_{n}^{-R} \;+\; T_{n}^{-R}\,\nu_{n}\,\kappa_{n}^{-R}\,\nu_{n}^{*}\,T_{n}^{-R}\Bigr) W^{*}
   \;-\; z^{n}\,\kappa_{n}^{-R}\,\nu_{n}^{*}\,T_{n}^{-R}\\
   \;& \quad -\; Z\,T_{n}^{-R}\,\nu_{n}\,\kappa_{n}^{-R}\,w^{-n}   
   \;+\; z^{n}\,\kappa_{n}^{-R} w^{-n}\, \\[6pt]
&= Z\,T_{n}^{-R}\,W^{*}
   \;+\;\bigl(z^{n}\,\kappa_{n}^{-R/2}\;-\;\varphi_{n}^{R}(z))\,
    \bigl(w^{-n}\,\kappa_{n}^{-R/2}\;-\;\varphi_{n}^R(w)^{\dag}) \\
&\quad -\;w^{-n}\,\bigl(z^{n}\,\kappa_{n}^{-R/2}\;-\;\varphi_{n}^{R}(z)\bigr)\,\kappa_{n}^{-R/2}
   \;-\;z^{n}\,\kappa_{n}^{-R/2}\,\bigl(w^{-n}\,\kappa_{n}^{-R/2}\;-\;\varphi_{n}^{R}(w)^{\dag}) \\
&\quad +\;z^{n}\,w^{-n}\,\kappa_{n}^{-R} \\[6pt]
&= K_{n-1}^{R}(w,z) + \varphi_{n}^{R}(z) \varphi_{n}^{R}(w)^{\dag}.
\end{align*}
Summing, the telescoping sum gives 1. as required.\\

\noindent For 3.
$$
F_{n}^{L}(z)=\begin{pmatrix}
\varphi_{n}^{L}(z) \\
\varphi_{n}^{R,*}(z) 
\end{pmatrix}, ~~ 
J=\begin{pmatrix}
I & 0 \\
0 & -I
\end{pmatrix}, ~~
\tilde{J}=\begin{pmatrix}
\bar{w}zI & 0 \\
0 & -I
\end{pmatrix}.
$$
Then 
$$
F_{n+1}^{L}(z)= A^{L}(\alpha_{n},z)F_{n}^{L}(z)
$$ 
and 
$$
A^{L}(\alpha_{n},w)^{\dag} J A^{L}(\alpha_{n},z)=\begin{pmatrix}
\bar{w}zI & 0 \\
0 & -I
\end{pmatrix} = \tilde{J}.
$$ 
Thus 
$$
F_{n+1}^{L}(w)^{\dag} J F_{n+1}^{L}(z) = F_{n}^{L}(w)^{\dag} A^{L}(\alpha_{n},w)^{\dag} J A^{L}(\alpha_{n},z) F_{n}^{L}(z) = F_{n}^{L}(w)^{\dag} \tilde{J} F_{n}^{L}(z)
$$ 
and hence
$$
\varphi_{n+1}^{L}(w)^{\dag}\varphi_{n+1}^{L}(z) - \varphi_{n+1}^{R,*}(w)^{\dag}\varphi_{n+1}^{R,*}(z)=\bar{w}z\varphi_{n}^{L}(w)^{\dag}\varphi_{n}^{L}(z) - \varphi_{n}^{R,*}(w)^{\dag}\varphi_{n}^{R,*}(z),
$$ 
which gives the second part of (a). Now, if we let 
$$
Q_{n}^{L}(z,w)=\varphi_{n+1}^{R,*}(w)^{\dag}\varphi_{n+1}^{R,*}(z) - \varphi_{n+1}^{L}(z)^{\dag}\varphi_{n+1}^{L}(z),
$$ 
we see 
\begin{eqnarray*}
Q_{n+1}^{L}(z,w) - Q_{n}^{L}(z,w) &=& \varphi_{n}^{R,*}(w)^{\dag}\varphi_{n}^{R,*}(z) - \bar{w}z\varphi_{n}^{L}(w)^{\dag}\varphi_{n}^{L}(z) \\
& &- \varphi_{n}^{R,*}(w)^{\dag}\varphi_{n}^{R,*}(z) + \varphi_{n}^{L}(w)^{\dag}\varphi_{n}^{L}(z)\\
&=&(1-\bar{w}z)\varphi_{n}^{L}(w)^{\dag}\varphi_{n}^{L}(z).
\end{eqnarray*}
Summing over $n$ the proof for (\ref{cd_right}) follows since $Q_{-1}^{L}(z,w)=0$.\\
The proof for (\ref{cd_left}) is similar.
\end{proof}

\end{document}
